\section{Discussion}\label{sec:discussion}
\subsection{What the semantic result establishes}
The method registers report meanings before grouping events. Field definitions and the common vocabulary are supplied, while missing permutations can be learned from provisional groups. New concepts, many-to-one mappings, and hierarchies require additional models.

Relative compatibility and absolute names remain distinct: common renaming preserves dot products and event groups. The known reference fixes names for hidden-entry evaluation, while event F1 evaluates grouped reports. Neither metric replaces the other.

Pair tables provide exact conditional optima for small vocabularies and avoid coordinate traps, but do not establish general superiority. Exact-minus-spectral intervals include zero, and forty scenes also allow reference-only recovery. Event groups remain approximate. Sequential and alternating maps coincide in all 150 nominal and same-class validation settings, supporting recovery from physical groups without benefit from repeated event updates.

Complete registration still leaves dense same-class F1 near 0.45, while shared wrong predictions can preserve event compatibility. Mapping accuracy, class correctness, and event F1 therefore measure different properties and need separate evaluation.

\subsection{Feasibility, accuracy, and dependence}
Feasibility is constructed, not evidence of correct identity. Physical-only constraints lose 0.0332 F1 against assignment with interval $[-0.0439,-0.0225]$, whereas boosted constraints gain 0.0108 with $[0.0020,0.0196]$. These scorer-dependent costs motivate reporting F1, ARI, $B^3$ F1, exact recovery, and violations together.

Concentration does not measure correctness. Unweighted scoring improves the nominal reference, whereas imbalance and receiver-specific errors reduce semantic gains. Event-shared errors can preserve agreement. Shared receiver errors also make $J$ composite, so pair calibration provides no calibrated group posterior. Replication intervals describe generator variation rather than group confidence.

The Gaussian bound requires a correct zero-mean error model and can fail under excess delay. Fragmentation conflicts with exclusivity; correlated false detections can form feasible wrong groups. Consolidation must identify repeated detections and propagate uncertainty. Shared artifacts require joint false-alarm modeling or nuisance-event evidence because lowering confidence alone cannot identify them. These unfitted models require measured detections and independent event labels. Pairwise admission still does not prove a common source position.

\subsection{Measured evidence and remaining validation}
LoED's conservative exclusions select a restricted traffic population. Equal payloads heard by complementary gateways can still arise from different transmissions, with no independent instrument resolving identity. The high F1 is conditional on this population and supervised preprocessing.

Only two of the original nine gateways remain eligible. Decoded-SF swaps supply no measured posteriors or higher-order semantic validation, while channel-discordant equal-payload pairs remain ambiguous. Training-metadata clock fitting removes one source of supervision, but population selection, scoring, and gate calibration still use references.

Field validation needs synchronized recordings from at least three receivers, independent event labels, documented native vocabularies, and measured class posteriors. WiSig addresses another identity task \cite{wisig}; the Sorbonne resource is a potential future source \cite{sorbonne2025,syed2023sniffer}, but suitable annotations remain to be established. No unverified substitute is evaluated.
